% SPDX-License-Identifier: Apache-2.0
% covers: LiteBridge
\datasheet{LiteBridge}{The AXI4 to AXI-Lite bridge: one AXI4 link to any count of AXI-Lite peripherals, by an address map.}

\dsfacts{\code{lib/parts/src/bus/axi\_lite.rs}}%
{\code{txhdl\_parts::bus::axi\_lite::LiteBridge<N, M, ..>}, with \code{M} an address map, \code{txhdl::map::AddrMap<N>}}%
{\code{//docs:axi}, \code{//docs:vreteno}, \code{//docs:eth}, \code{//docs:hdmi}}

\subsection*{Function}

AXI-Lite is AXI4 with no identifier, no burst and no \code{last}, so
one transaction is one beat on each channel it uses. A bridge stands
between an AXI4 link and that many AXI-Lite peripherals. On its AXI4
side it takes the \code{aw}, \code{ar} and \code{w} an \code{AxiPer}
would take, and gives the \code{b} and \code{r} it would give. So it
goes where an \code{AxiPer} would: behind a host's tracker, or behind
one of a router's peripheral ports.

Its peripheral side is arrays of \code{N} ports: it drives
\code{aws}, \code{ars} and \code{ws}, and reads \code{bs} and
\code{rs}, entry \textit{i} of each being peripheral \textit{i}'s, so
the netlist's ports are \code{aws\_}\textit{i} and the like. An
address phase there is a \code{LiteAddr}: \code{addr} and
\code{prot}. A write word is a \code{LiteW}: \code{data} and
\code{strb}. \code{axi\_lite} makes the five channels of such a link.
One unit serves every count; the lowering unrolls its loops over the
arrays when it runs (issue 500). Until issue 647 a macro,
\code{lite\_bridge!}, wrote \code{LiteBridge1} to \code{LiteBridge8}
instead, one unit per count, with each range a pair of const
parameters.

\subsection*{Parameters}

\begin{center}\footnotesize
\begin{tabular}{@{}lp{0.7\textwidth}@{}}
\toprule
Parameter & Meaning \\
\midrule
\code{A} & The address width, in bits, on both sides. \\
\code{D} & The data width, in bits, on both sides. \\
\code{S} & The write strobe width, one bit per byte lane, so \code{D / 8}. \\
\code{I} & The identifier width of the AXI4 side, in bits. \\
\code{N} & The count of peripherals, one or more. \\
\code{M} & The address map: a type implementing \code{AddrMap<N>},
whose \code{RANGES} holds each peripheral's base and mask, in the
order of the ports. \\
\bottomrule
\end{tabular}
\end{center}

\subsection*{Address map}

An address \code{a} is peripheral \textit{i}'s when
\code{(a \& mask) == base} for \code{(base, mask) = M::RANGES[}\textit{i}\code{]}.
The bridge decodes a burst's first address, and the lowest range that
matches wins. An address in no range is a hole. The tables below are
generated for \code{LiteBridge<1, SerialMap, 32, 32, 4, 2>}, the bridge
in front of Vreteno's serial port, whose map is this one:

\begin{center}\footnotesize
\begin{tabular}{@{}llll@{}}
\toprule
Peripheral & \code{BASE} & \code{MASK} & In Vreteno \\
\midrule
0 & \code{0x3000} & \code{0xf000} & The serial port \\
\bottomrule
\end{tabular}
\end{center}

The board's design uses a bridge of six, \code{SlotMap}, at
\code{0x3000} to \code{0x3500} with mask \code{0xffff\_ff00} each, and
bridges of one at \code{0x0c00\_0000} (mask \code{0xfc00\_0000}, in
front of its interrupt controller) and at \code{0x1000\_0000} (mask
\code{0xffff\_0000}, the debug module). \code{ex\_axi\_lite} uses a
bridge of two with ranges at \code{0x1000} and \code{0x2000}, mask
\code{0xf000} each.

\dstables{LiteBridge}

\subsection*{Behaviour}

\begin{itemize}
\item \textbf{Taking a burst.} It takes a burst only when none is in
progress (\code{busy} clear). When both address channels offer,
\code{rfirst} says which goes first. After a read the write channel
goes first, and after a write the read channel.
\item Taking a burst records its identifier in \code{xid}, its address
in \code{addr}, its \code{len} in \code{left}, its \code{prot}, and
its peripheral in \code{sel}, a bit each. A hole has no bit set.
\code{sel} is \code{N} bits wide; a bridge of one has a one-bit
\code{sel}, which the netlist writes as a single bit.
\item \code{stride} is how far the address moves per beat:
\code{1 << size} bytes, or zero for a \code{Fixed} burst. Every other
burst kind moves as \code{Incr} does.
\item \textbf{A write beat.} The bridge takes the burst's next beat
from \code{w} when the chosen peripheral has room on both its
\code{aw} and its \code{w}. It sends that peripheral \code{addr} and
\code{prot} on \code{aw}, and the beat's \code{data} and \code{strb}
on \code{w}. To a hole, the beat is taken and dropped.
\item \textbf{A read beat.} The bridge sends \code{addr} and
\code{prot} on the chosen peripheral's \code{ar} when it has room.
\item Once a beat's request is out (\code{sent}), the bridge waits for
that peripheral's answer. A hole's answer is there at once:
\code{DecErr}, with zero data for a read.
\item \textbf{A read's answer} goes up on \code{r} as a beat of the
burst, when \code{r} has room. It has the burst's identifier, the word,
the peripheral's response, and \code{last} on the final beat.
\item \textbf{A write's answer} is taken when \code{b} has room, or when
more beats are to come. \code{wresp} keeps the first response that is
not \code{Okay}. After the last beat one \code{B} goes up, with the
burst's identifier and that response.
\item After each answer, \code{addr} moves on by \code{stride} and
\code{left} counts down. The answer to the last beat clears
\code{busy}.
\end{itemize}

\subsection*{Verification}

\begin{itemize}
\item \code{bursts\_cross\_the\_bridge\_a\_beat\_at\_a\_time} in
\code{bus::axi\_lite}, in \code{//lib/parts:parts\_test}, puts
\code{LiteBridge<2, TwoMap, 16, 32, 4, 2>} behind
a host tracker, with two simulated memories. It writes three words and
reads four, reads one word three times in a fixed burst, uses the
second peripheral, and sends a write and a read of two beats to a hole.
It checks every word and response, and counts ten transactions on the
first memory and two on the second.
\item \code{ex\_axi\_lite} runs the same kind of bursts against two
banks of registers written as hardware. The build simulates the
bridge's netlist against that run under nvc and Verilator:
\begin{itemize}
\item \code{//docs:axi\_lite\_sim\_axi\_lite\_bridge\_axi\_lite\_bridge\_tb\_test};
\item \code{//docs:axi\_lite\_vsim\_axi\_lite\_bridge\_test}.
\end{itemize}
\item Vreteno's run lowers its serial port's bridge, the
bridge of one of the tables, and the build simulates that netlist
against the run: \code{//docs:vreteno\_sim\_ubridge\_ubridge\_tb\_test}
and \code{//docs:vreteno\_vsim\_ubridge\_test}.
\item \code{//cpu/vreteno:lockstep\_test} and \code{//soc:demo\_test} run
a bridge of one in simulation.
\end{itemize}

\subsection*{Used by}

A bridge of two in the example \code{ex\_axi\_lite}. A bridge of one
in front of the serial port in Vreteno's run
(\code{cpu/vreteno/src/run.rs} and \code{cpu/vreteno/src/main.rs}) and
in \code{soc/src/lib.rs}; in the board's design, a bridge of six as its
field \code{puart} and bridges of one as \code{pplic} and \code{pdm}.
A bridge of one also puts each peripheral of the examples behind a
host's tracker: \code{ex\_eth}, \code{ex\_hdmi}, \code{ex\_sd} and the
rest.

\subsection*{Limits}

It serves one burst at a time, for all its peripherals together, since
AXI-Lite has no identifier to tell two apart. A burst's peripheral is
chosen on its first address, and every beat of the burst goes there.
It treats a \code{Wrap} burst as \code{Incr}. It counts a write's beats
from \code{len} and does not read \code{last} on \code{w}. Of the
address phase only \code{addr} and \code{prot} reach a peripheral;
\code{size} sets the stride, and \code{lock}, \code{cache},
\code{qos} and \code{region} are dropped.
